
% =============================================================================
% Project : THE PHILOSOPHY OF VIABILITY
% Book    : Book I — Ontology of Viability
% Book ID : VS.B01
% Part    : Part I — Identity and Manifestation
% Chapter : Chapter 08 — Comparative Philosophical Examination
% Chapter ID : VS.B01.P01.C08
%
% File :
%   4_ASSETS/tables/
%   VS.B01.P01.C08.S09.SS14.P003.TL01.hume_claims_and_comparative_limits.tex
%
% Section / Chapter / Asset:
%   Table TL01 — Hume's Claims and Their Comparative Limits
%
% Structural Identifier:
%   VS.B01.P01.C08.S09.SS14.P003.TL01
%
% Asset Type : TL
%
% Purpose:
%   To distinguish Hume's principal arguments according to
%   their subject matter, explanatory function, and limits.
%
% Representational Limit:
%   The table is an analytical summary. It does not imply that
%   Hume's arguments form one completed or internally
%   uncontested ontological system.
%
% =============================================================================

\begin{table}[htbp]
\centering
\small

\caption{Hume's Principal Claims and Their Comparative Limits}
\label{tab:VS.B01.P01.C08.S09.SS14.P003.TL01}

\begin{tabularx}{\textwidth}{
    >{\raggedright\arraybackslash}p{2.25cm}
    >{\raggedright\arraybackslash}X
    >{\raggedright\arraybackslash}X
    >{\raggedright\arraybackslash}X
}
\toprule
\textbf{Problem}
&
\textbf{Hume's account}
&
\textbf{What it explains}
&
\textbf{What it does not establish}
\\
\midrule

Copy Principle
&
Simple ideas are derived from and resemble simple impressions,
subject to the qualifications raised by Hume himself.
&
The origin and intelligibility of ideas.
&
That whatever cannot be traced to an impression does not exist.
\\

Induction
&
Custom or habit explains the transition from observed
regularities to expectations about further instances.
&
The formation of inductive belief and expectation.
&
A deductive proof that the future must resemble the past,
or that every inductive belief is false.
\\

Causation
&
Constant conjunction and customary transition explain
the idea of necessary connection.
&
The formation of causal expectation and the idea of necessity.
&
An uncontested conclusion that causal relations exist
only in the mind or that objective causal relations
have been fully explained.
\\

Identity of objects
&
Relations among successive parts facilitate the imagination's
attribution of identity to changing or interrupted objects.
&
Why certain successions are ordinarily treated as one object.
&
That the ontological identity of an object is identical
with the psychological attribution of sameness.
\\

Personal identity
&
Resemblance, contiguity, causation, memory, and imagination
contribute to the attribution of identity to successive
perceptions.
&
The structure of Hume's account of personal identity.
&
A complete and uncontested explanation of the subject
or unity of experience.
\\

The Appendix
&
Hume acknowledges a difficulty in explaining the principle
that connects distinct perceptions.
&
The limits and unresolved tension within the earlier account.
&
That the bundle interpretation is a fully settled
and unproblematic final doctrine.
\\

\bottomrule
\end{tabularx}

\caption*{\footnotesize
The table distinguishes the explanatory task performed by
each argument from the broader conclusions that might be
drawn from it. The latter require additional argument and
must not be treated as automatic consequences of Hume's
textual claims.
}

\end{table}

% =============================================================================
% END OF TABLE TL01
% =============================================================================
